\subsection[class \texttt{Model}]{\classpymodel}\label{sec:class_py_model}
The \gretlpy module implements the \classpymodel class,
which is used to manage the definition, estimation,
and output of models. The constructor of this class
has the following signature:
\begin{lstlisting}[style=pythonstyle]
def __init__(self, name, *args, **kwargs)
\end{lstlisting}
Here, \pythcode{name} specifies the estimator's name,
and \pythcode{*args} and \pythcode{**kwargs} represent
additional (including optional) arguments, as detailed
in Section~\ref{sec:estimators_python}, with examples
provided in Section~\ref{sec:estimators_python_examples}.

\vspace*{20pt}\noindent The \classpymodel class provides the following public methods:

\gretlfpycmd{aic()}
\estdesc{Returns value of Akaike criterion.}
\begin{lstlisting}[style=pythonstyle]
float = aic(self)
\end{lstlisting}

\gretlfpycmd{bic()}
\estdesc{Returns value of Schwartz criterion.}
\begin{lstlisting}[style=pythonstyle]
float = bic(self)
\end{lstlisting}

\gretlfpycmd{coeff()}
\estdesc{Returns coefficient estimates.}
\begin{lstlisting}[style=pythonstyle]
numpy.ndarray = coeff(self)
\end{lstlisting}

\gretlfpycmd{fit()}
\estdesc{Executes the underlying implementation of
  the specified estimator in \gretlpy.}
\begin{lstlisting}[style=pythonstyle]
gretl.Model = fit(self)
\end{lstlisting}

\gretlfpycmd{hqc()}
\estdesc{Returns value of Hannan-Quinn criterion.}
\begin{lstlisting}[style=pythonstyle]
float = hqc(self)
\end{lstlisting}

\gretlfpycmd{id()}
\estdesc{Model's ID (assigned by \lgretl).}
\begin{lstlisting}[style=pythonstyle]
int = id(self)
\end{lstlisting}

\gretlfpycmd{lnl()}
\estdesc{Returns the log-likelihood for the last estimated model.}
\begin{lstlisting}[style=pythonstyle]
int = lnl(self)
\end{lstlisting}

\gretlfpycmd{oddsratios()}\index{logit()}
\estdesc{Returns oddsratios for logit model (if available).}
\begin{lstlisting}[style=pythonstyle]
numpy.ndarray = oddsratios(self, show=bool)
\end{lstlisting}
\generaldesc{This methods retrieves and pretty prints
	oddsratios for \gretlcmd{logit} model:
	\begin{itemize}
		\item \pythcode{show=bool}: If set to \pythcode{True},
			table with oddsratios is printed (default:~\pythcode{True}).
\end{itemize}}

\gretlfpycmd{parnames()}
\estdesc{Returns names of coefficients.}
\begin{lstlisting}[style=pythonstyle]
str = parnames(self)
\end{lstlisting}

\gretlfpycmd{rsq()}
\estdesc{Returns R-squared.}
\begin{lstlisting}[style=pythonstyle]
float = rsq(self)
\end{lstlisting}

\gretlfpycmd{sigma()}
\estdesc{Returns standard errors of residuals.}
\begin{lstlisting}[style=pythonstyle]
float = sigma(self)
\end{lstlisting}

\gretlfpycmd{stderr()}
\estdesc{Returns standard errors of coefficient estimates.}
\begin{lstlisting}[style=pythonstyle]
numpy.ndarray = stderr(self)
\end{lstlisting}

\gretlfpycmd{uhat()}
\estdesc{Returns residuals.}
\begin{lstlisting}[style=pythonstyle]
list = uhat(self)
\end{lstlisting}
\generaldesc{Note: in case of \gretlcmd{biprobit}\index{biprobit()},
  \texttt{bp\_uhat} entry of \gretlcmd{\$model} bundle is returned.}

\gretlfpycmd{yhat()}
\estdesc{Returns fitted values.}
\begin{lstlisting}[style=pythonstyle]
list = yhat(self)
\end{lstlisting}
\generaldesc{Note: in case of \gretlcmd{biprobit}\index{biprobit()},
  \texttt{bp\_yhat} entry of \gretlcmd{\$model} bundle is returned.}

\vspace*{20pt}\noindent The \classpymodel class provides the following public attributes:

\gretlfpycmd{args}
\estdesc{Options passed as \pythcode{*args}.}
\begin{lstlisting}[style=pythonstyle]
tuple = args
\end{lstlisting}

\gretlfpycmd{computed}
\estdesc{Indicated if the model was estimated.}
\begin{lstlisting}[style=pythonstyle]
bool = computed
\end{lstlisting}

\gretlfpycmd{options}
\estdesc{Options passed as \pythcode{**kwargs}.}
\begin{lstlisting}[style=pythonstyle]
dict = options
\end{lstlisting}

\gretlfpycmd{results}
\estdesc{Contents of the \lgretl's \gretlcmd{\$model} bundle
  \cite[see][]{gretl_command} for details.}
\begin{lstlisting}[style=pythonstyle]
dict = results
\end{lstlisting}

\gretlfpycmd{summary}
\estdesc{Model's output.}
\begin{lstlisting}[style=pythonstyle]
str = summary
\end{lstlisting}
